\subsection{Measures and bounded continuous functions}

Recall (GT, IX, §1, No.~5, Def. 4) that a topological space T is said to be completely regular if it is uniformizable and Hausdorff. This is equivalent to saying (loc. cit., Prop. 3) that T is homeomorphic to a subspace of a compact space. If T is completely regular, then every positive lower semi-continuous function f on T is the upper envelope of the increasing directed set of elements of \(\mathcal{C}_{+}^{b}(T)\) that are \(\leqslant f\) , and every positive and bounded upper semi-continuous function g is the lower envelope of the decreasing directed set of elements of \(\mathcal{C}_{+}^{b}(T)\) that are \(\geqslant g\) (loc. cit., §1, No.~6, Prop. 5). We shall need the following lemma:

Lemma. --- Let T be a completely regular space, K a compact subset of T, and U an open subset of T containing K.

\begin{enumerate}
\def\labelenumi{\alph{enumi})}
\item
  There exists an open subset \(\mathbf{U}'\) of \(\mathbf{T}\) such that \(\mathbf{K} \subset \mathbf{U}' \subset \overline{\mathbf{U}'} \subset \mathbf{U}\) .
\item
  Let \(f\) be a continuous function defined on \(K\) with values in an interval \(I\) of \(\mathbf{R}\) (resp. in \(\mathbf{C}\) ). There exists a bounded continuous function \(f'\) on \(T\) , with values in \(I\) (resp. in \(\mathbf{C}\) ), that extends \(f\) and is zero on \(T - U\) .
\end{enumerate}

It suffices to treat the case that T is a subspace of a compact space X. Let V be an open subset of X such that \(V \cap T = U\) ; denote by \(V'\) an open set in X containing K such that \(\overline{V'} \subset V\) , by g a continuous function on X with values in I (resp. in C) extending f and zero on X - V (GT, IX, §4, No.~1, Prop. 1). The condition a) is satisfied by taking \(U' = V' \cap T\) , and b) by taking \(f'\) to be the restriction of g to T.

PROPOSITION 1. --- Let T be a completely regular space.

\begin{enumerate}
\def\labelenumi{\alph{enumi})}
\tightlist
\item
  Let \(\mu\) be a positive measure on \(\mathbf{T}\) , and \(f\) a numerical function \(\geqslant 0\) defined on \(\mathbf{T}\) and lower semi-continuous (resp. upper semi-continuous, finite, with compact support). Then
\end{enumerate}

\[
\mu^ {\bullet} (f) = \sup _ {g \in I _ {f}} \mu^ {\bullet} (g) \quad (\text { resp. } \mu^ {\bullet} (f) = \inf _ {g \in S _ {f}} \mu (g)),\tag{1}
\]

where \(I_{f}\) (resp. \(S_{f}\) ) denotes the set of bounded continuous functions g such that \(0 \leqslant g \leqslant f\) (resp. \(g \geqslant f\) ).

\begin{enumerate}
\def\labelenumi{\alph{enumi})}
\setcounter{enumi}{1}
\tightlist
\item
  Let \(\theta\) be a complex measure on \(\mathbf{T}\) , and \(f\) a numerical function \(\geqslant 0\) defined on \(\mathbf{T}\) and lower semi-continuous. Then
\end{enumerate}

\[
| \theta | ^ {\bullet} (f) = \sup _ {g} | \theta (g) |,\tag{2}
\]

where \(g\) runs over the set of bounded and \(|\theta|\) -integrable continuous complex functions such that \(|g| \leqslant f\) .

The first of the formulas (1) is obvious, because \(I_{f}\) is an increasing directed set of continuous functions whose upper envelope is f, and one can apply Prop. 5 of §1, No.~6. The same proposition will imply the second formula, if we show that \(S_{f}\) contains a \(\mu\) -integrable bounded continuous function. Thus, let K be the support of f, and M the supremum of f; since K is compact, M is finite (GT, IV, §6, No.~2, Th. 3). Let U be an open set containing K and such that \(\mu^{\bullet}(U) < +\infty\) ; there exists (Lemma) a continuous function g with values in {[}0, M{]}, equal to M on K and zero outside U; then \(g \in S_{f}\) and \(\mu^{\bullet}(g) \leqslant M\mu^{\bullet}(U) < +\infty\) .

Let us pass to \(b\) . It clearly suffices to show that \(|\theta|^{\bullet}(f) \leqslant \sup |\theta(g)|\) .

Let \(a\) and \(b\) be two real numbers such that \(a < b < |\theta|^{\bullet}(f)\) . By (1), there exists a function \(h \in \mathcal{C}_{+}^{b}(\mathrm{T})\) such that \(h \leqslant f\) and \(|\theta|^{\bullet}(h) > b\) ; denote by M the supremum of \(h\) . By the definition of \(|\theta|^{\bullet}\) (§1, No.~2, Def. 4), there exists a compact subset K of T such that \(|\theta|_{\mathrm{K}}^{\bullet}(h_{\mathrm{K}}) > b\) . There then exists a continuous complex function \(j\) on K such that \(|j| \leqslant h_{\mathrm{K}}\) and \(|\theta_{\mathrm{K}}(j)| > b\) (Ch. III, §1, No.~6). Let us choose an open set U containing K and such that \(|\theta|^{\bullet}(\mathrm{U} - \mathrm{K}) \leqslant \frac{b - a}{\mathrm{M}}\) (§1, No.~9, Props. 13 and 14); extend \(j\) to a continuous complex function \(k\) on T, zero outside U (Lemma); for every \(t \in \mathrm{T}\) , set

\[
g (t) = \left\{ \begin{array}{l l} k (t) & \text { if } | k (t) | \leqslant h (t) \\ \frac {k (t)}{| k (t) |} h (t) & \text { if } | k (t) | > h (t). \end{array} \right.\tag{3}
\]

Clearly \(|g| \leqslant h \leqslant f\) , and \(g = j\) on \(\mathbf{K}\) , therefore \(\left|\theta_{\mathrm{K}}(j) - |\theta(g)|\right| = \left|\left|\theta(j^0)\right| - |\theta(g)|\right| \leqslant |\theta|^{\bullet}(|j^0 - g|) \leqslant \mathbf{M} \cdot |\theta|^{\bullet}(\mathbf{U} - \mathbf{K}) \leqslant b - a\) , consequently \(|\theta(g)| > a\) . Let us show on the other hand that \(g\) is a continuous function:

since \(a\) is subject to the sole condition \(a < |\theta|^{\bullet}(f)\) , this will imply that the second member of (2) is \(\geqslant\) the first, whence the proposition. Now, let \(F\) (resp. \(F'\) ) be the set of \(t \in T\) such that \(|k(t)| \leqslant h(t)\) (resp. \(|k(t)| \geqslant h(t)\) ). These sets being closed, and their union being \(T\) , it will suffice to show that \(g_{F}\) and \(g_{F'}\) are continuous: now, this property is obvious for \(g_{F} = k_{F}\) , and it is so for \(g_{F'}\) at the points where \(k(t) \neq 0\) ; on the other hand, if \(t \in F'\) is such that \(k(t) = 0\) , then also \(h(t) = 0\) , and the inequality \(|g| \leqslant h\) implies that \(g\) is continuous at the point \(t\) .

Remarks. --- 1) Let \(f\) be a positive lower semi-continuous function, and let \(J_f\) be the set of positive bounded continuous functions zero outside a \(\mu\) -integrable open set and bounded above by \(f\) . One can show that \(f\) is the upper envelope of \(J_f\) and that \(\mu^\bullet(f) = \sup_{g \in L} \mu(g)\) .

\begin{enumerate}
\def\labelenumi{\arabic{enumi})}
\setcounter{enumi}{1}
\tightlist
\item
  If the measure \(\mu\) is bounded, the formula \(\mu^{\bullet}(f) = \inf_{g\in \mathbb{S}_f}\mu (g)\) is obviously
\end{enumerate}

valid for every function f that is upper semi-continuous, positive and bounded.

PROPOSITION 2. --- Let \(\eta\) and \(\eta'\) be two complex measures on a completely regular space \(T\) , such that \(\eta(f) = \eta'(f)\) for every function \(f \in \mathcal{C}^b(T)\) that is integrable for \(|\eta|\) and \(|\eta'|\) . Then \(\eta = \eta'\) .

Let us take up again the proof of the second part of Proposition 1, on setting \(\theta = \eta - \eta'\) . We can require the open set U to be integrable for \(|\eta|\) and \(|\eta'|\) . The function \(g\) is then integrable for these two measures, and the relation \(\theta(g) = 0\) implies \(a < 0\) ; therefore \(|\theta|^{\bullet}(f) = 0\) for every positive lower semi-continuous function \(f\) , whence finally \(|\theta| = 0\) , on taking \(f = +\infty\) .

PROPOSITION 3. --- Let \(\mu\) be a positive measure on a completely regular space \(T\) , and let \(p \in [1, +\infty[\) . The space \(\mathcal{H}\) of functions \(f \in \mathcal{C}^b(T)\) , whose support is contained in a \(\mu\) -integrable open set, is dense in \(\mathcal{L}^p(\mu)\) .

By Prop. 15 of §1, No.~10, it suffices to show that if K is compact in T, and if g is the extension to T by 0 of a function in \(\mathcal{C}_{+}(K)\) between 0 and 1, then there exists a function \(f \in \mathcal{C}_{+}^{b}(T)\) , with support contained in a \(\mu\) -integrable open set, such that \(\|f - g\|_{p}\) is arbitrarily small. Now, let \(\varepsilon\) be a number \textgreater0, U an open neighborhood of K such that \(\mu^{\bullet}(U - K) < \varepsilon\) , V an open neighborhood of K such that \(\overline{V} \subset U\) , and f a function with values in {[}0,1{]}, continuous, equal to g on K and to 0 outside V (Lemma). The function \(|f - g|^{p}\) is then bounded above by \(\varphi_{U-K}\) ; therefore \(\|f - g\|_{p} \leqslant \varepsilon^{1/p}\) , which establishes the proposition.

Remark 3). --- There is an analogous statement for functions with values in a Banach space F: the subspace \(\mathcal{H} \otimes \mathrm{F}\) of \(\mathcal{C}^b (\mathrm{T};\mathrm{F})\) is dense in \(\mathcal{L}_{\mathrm{F}}^{p}(\mu)\) .

PROPOSITION 4. --- In order that a bounded complex measure \(\theta\) on a completely regular space \(T\) be positive, it is necessary and sufficient that \(\theta(f) \geqslant 0\) for every function \(f \in \mathcal{C}_{+}^{b}(T)\) .

Necessity is obvious. To establish sufficiency, let us take up again the proof of the preceding proposition, on taking p = 1 and \(\mu = |\theta|\) ; the notations being the same, the relation \(\mu^{\bullet}(|f - g|) \leqslant \varepsilon\) and the inequality \(\theta(f) \geqslant 0\) imply \(\theta_{\mathrm{K}}(g_{\mathrm{K}}) = \theta(g) \geqslant -\varepsilon\) ; since \(g_{K}\) is an arbitrary element of \(\mathcal{C}(K)\) between 0 and 1, the measure \(\theta_{K}\) is positive; the compact set K being arbitrary, this means that \(\theta\) is positive.

